\documentclass[a4paper,12pt]{article}

\usepackage[margin=1in]{geometry}
\usepackage{listings}
\usepackage{fancyhdr}

% ==========================================
% PAGE STYLE
% ==========================================

\pagestyle{fancy}

\fancyhead{}

\fancyfoot{}

\fancyfoot[C]{\thepage}

\renewcommand{\headrulewidth}{0pt}

\renewcommand{\footrulewidth}{0.4pt}

% ==========================================
% JAVA CODE STYLE
% ==========================================

\lstdefinestyle{javacode}{
    language=Java,
    basicstyle=\ttfamily\small,
    numbers=left,
    numberstyle=\tiny,
    numbersep=10pt,
    frame=single,
    breaklines=true,
    showstringspaces=false,
    columns=fullflexible,
    keepspaces=true,
    tabsize=4
}

% ==========================================
% OUTPUT STYLE
% ==========================================

\lstdefinestyle{outputstyle}{
    basicstyle=\ttfamily\small,
    frame=single,
    breaklines=true,
    showstringspaces=false,
    columns=fullflexible,
    keepspaces=true
}

% ==========================================
% EXPERIMENT COUNTER
% ==========================================

\newcounter{experiment}
\setcounter{experiment}{55}

% ==========================================
% START EXPERIMENT
%
% #1 = Date
% #2 = Aim
% #3 = Java file
% ==========================================

\newcommand{\startexperiment}[3]{

\stepcounter{experiment}

\begin{center}

    {\LARGE\textbf{JAVA PROGRAMMING LAB}}

\end{center}

\vspace{0.8cm}

\noindent

\textbf{Experiment No.: \theexperiment}

\hfill

\textbf{Date: #1}

\vspace{0.2cm}

\hrule

\vspace{0.8cm}

\noindent

\textbf{Aim:}

\vspace{0.3cm}

#2

\vspace{0.7cm}

\noindent

\textbf{Program:}

\vspace{0.3cm}

\lstinputlisting[style=javacode]{#3}

\vspace{0.7cm}

\noindent

\textbf{Output:}

\vspace{0.3cm}

}

% ==========================================
% END EXPERIMENT
% ==========================================

\newcommand{\finishExperiment}{
    \newpage
}

% ==========================================
% DOCUMENT
% ==========================================

\begin{document}
\setcounter{page}{73}

% ==========================================
% PROGRAM 56
% ==========================================

\startexperiment{13/08/2026}{WAP to accept and display employee details such as employee ID, name and salary.}{56.java}

\begin{lstlisting}[style=outputstyle]
Enter empid: 101
Enter empname: Raghu
Enter salary: 75000
Employee Details
Employee ID: 101
Employee Name: Raghu
Employee Salary: 75000.0
\end{lstlisting}

\finishExperiment

% ==========================================
% PROGRAM 57
% ==========================================

\startexperiment{13/08/2026}{WAP to accept the number of rows and columns of a matrix, read its elements and display its transpose.}{57.java}

\begin{lstlisting}[style=outputstyle]
Enter the number of rows: 3
Enter the number of columns: 3
Enter the values in 3 x 3 matrix:
1
2
3
4
5
6
7
8
9
Matrix is:
1 2 3 
4 5 6 
7 8 9 
Transpose of Matrix:
1 4 7 
2 5 8 
3 6 9 
\end{lstlisting}

\finishExperiment

% ==========================================
% PROGRAM 58
% ==========================================

\startexperiment{13/08/2026}{WAP to display numbers from 10 to 100 with a difference of 10 using a loop.}{58.java}

\begin{lstlisting}[style=outputstyle]
10
20
30
40
50
60
70
80
90
100
\end{lstlisting}

\finishExperiment

% ==========================================
% PROGRAM 59
% ==========================================

\startexperiment{13/08/2026}{WAP to find and display the square of a given number using the Math.pow() function.}{59.java}

\begin{lstlisting}[style=outputstyle]
Enter a number: 12
The square of 12 is: 144
\end{lstlisting}

\finishExperiment

% ==========================================
% PROGRAM 60
% ==========================================

\startexperiment{13/08/2026}{WAP to find and display the sine and cosine values of a given number using Math.sin() and Math.cos() functions.}{60.java}

\begin{lstlisting}[style=outputstyle]
Enter angle in degrees: 90
Sine value = 1.0
Cosine value = 6.123233995736766E-17
\end{lstlisting}

\finishExperiment

% ==========================================
% PROGRAM 61
% ==========================================

\startexperiment{13/08/2026}{WAP to find and display the maximum value between two given values using the Math.max() function.}{61.java}

\begin{lstlisting}[style=outputstyle]
Enter first value: 22
Enter second value: 55
Maximum value = 55
\end{lstlisting}

\finishExperiment

% ==========================================
% PROGRAM 62
% ==========================================

\startexperiment{13/08/2026}{WAP to find and display the power of one number raised to another number using the Math.pow() function.}{62.java}

\begin{lstlisting}[style=outputstyle]
Enter base value: 4
Enter power value: 4
Power = 256.0
\end{lstlisting}

\finishExperiment

% ==========================================
% PROGRAM 63
% ==========================================

\startexperiment{13/08/2026}{WAP to find and display the smaller value between two given numbers using the Math.min() function.}{63.java}

\begin{lstlisting}[style=outputstyle]
Enter first value: 20
Enter second value: 50
Smaller value = 20
\end{lstlisting}

\finishExperiment

% ==========================================
% PROGRAM 64
% ==========================================

\startexperiment{13/08/2026}{WAP to reverse and display a given string.}{64.java}

\begin{lstlisting}[style=outputstyle]
Enter a string: Aditya
Reversed string = aytidA
\end{lstlisting}

\finishExperiment

% ==========================================
% PROGRAM 65
% ==========================================

\startexperiment{13/08/2026}{WAP to display a given string using string methods.}{65.java}

\begin{lstlisting}[style=outputstyle]
Enter a string: Aditya
String = Aditya
Length = 6
Uppercase = ADITYA
Lowercase = aditya
First character = A
\end{lstlisting}

\finishExperiment

% ==========================================
% PROGRAM 66
% ==========================================

\startexperiment{13/08/2026}{WAP to display a given string and demonstrate string manipulation.}{66.java}

\begin{lstlisting}[style=outputstyle]
Enter a string: Aditya Raj
Original String = Aditya Raj
Length = 10
Uppercase = ADITYA RAJ
Lowercase = aditya raj
First character = A
Substring = Adity
Replaced String = AdityA RAj
\end{lstlisting}

\finishExperiment

% ==========================================
% PROGRAM 67
% ==========================================

\startexperiment{13/08/2026}{WAP to replace a character in a given string with another character.}{67.java}

\begin{lstlisting}[style=outputstyle]
Enter a string: ranchi
Enter character to replace: r
Enter new character: R
Original String = ranchi
Replaced String = Ranchi
\end{lstlisting}

\finishExperiment

% ==========================================
% PROGRAM 68
% ==========================================

\startexperiment{13/08/2026}{WAP to find and display the length of a given string.}{68.java}

\begin{lstlisting}[style=outputstyle]
Enter a string: Aditya
Length of string = 6
\end{lstlisting}

\finishExperiment

% ==========================================
% PROGRAM 69
% ==========================================

\startexperiment{20/08/2026}{WAP accept any two number from the keyboard and display their sum.}{69.java}

\begin{lstlisting}[style=outputstyle]
Enter first number: 20
Enter second number: 50
Sum = 70
\end{lstlisting}

\finishExperiment

% ==========================================
% PROGRAM 70
% ==========================================

\startexperiment{20/08/2026}{WAP to accept and display a character from the keyboard. Using scanner class.}{70.java}

\begin{lstlisting}[style=outputstyle]
Enter a character: A
The entered character is: A
\end{lstlisting}

\finishExperiment

% ==========================================
% PROGRAM 71
% ==========================================

\startexperiment{20/08/2026}{WAP to accept and display a string or a name from the keyboard. Using scanner class.}{71.java}

\begin{lstlisting}[style=outputstyle]
Enter your name: Aditya
Your name is: Aditya
\end{lstlisting}

\finishExperiment

% ==========================================
% PROGRAM 72
% ==========================================

\startexperiment{20/08/2026}{WAP to accept and display an integer from the keyboard. Using scanner class}{72.java}

\begin{lstlisting}[style=outputstyle]
Enter an integer: 12
The entered integer is: 12
\end{lstlisting}

\finishExperiment

% ==========================================
% PROGRAM 73
% ==========================================

\startexperiment{20/08/2026}{WAP to accept and display a float number from the keyboard. Using scanner class.}{73.java}

\begin{lstlisting}[style=outputstyle]
Enter a float number: 22.34
The entered float number is: 22.34
\end{lstlisting}

\finishExperiment

% ==========================================
% PROGRAM 74
% ==========================================

\startexperiment{20/08/2026}{WAP to accept and display employee details. Using scanner class.}{74.java}

\begin{lstlisting}[style=outputstyle]
Enter empid: 101
Enter empname: Aditya
Enter salary: 100000
Employee Details
Employee ID: 101
Employee Name: Aditya
Employee Salary: 100000.0
\end{lstlisting}

\finishExperiment

% ==========================================
% PROGRAM 75
% ==========================================

\startexperiment{20/08/2026}{WAP to accept any integer from the keyboard and test if it is even or not.}{75.java}

\begin{lstlisting}[style=outputstyle]
Enter an integer: 24
24 is an even number.
\end{lstlisting}

\finishExperiment

% ==========================================
% PROGRAM 76
% ==========================================

\startexperiment{20/08/2026}{WAP to accept any two float number from the keyboard and display their sum.}{76.java}

\begin{lstlisting}[style=outputstyle]
Enter first float number: 33.45
Enter second float number: 65.54
Sum = 98.990005
\end{lstlisting}

\finishExperiment

% ==========================================
% PROGRAM 77
% ==========================================

\startexperiment{20/08/2026}{WAP to create a 1D array and read its elements by using a loop display them one by one.}{77.java}

\begin{lstlisting}[style=outputstyle]
Enter size of array: 3
Enter element 1: 1
Enter element 2: 2
Enter element 3: 3
Array elements are:
1
2
3
\end{lstlisting}

\finishExperiment

% ==========================================
% PROGRAM 78
% ==========================================

\startexperiment{20/08/2026}{WAP to accepts the marks of a students into 1D array from the keyboard and finds total marks and percentage.}{78.java}

\begin{lstlisting}[style=outputstyle]
Enter number of subjects: 4
Enter marks of subject 1: 88
Enter marks of subject 2: 67
Enter marks of subject 3: 98
Enter marks of subject 4: 59
Total Marks = 312
Percentage = 78.0%
\end{lstlisting}

\finishExperiment

% ==========================================
% PROGRAM 79
% ==========================================

\startexperiment{20/08/2026}{WAP to accept any integer from the keyboard and check whether it is positive, negative, or zero.}{79.java}

\begin{lstlisting}[style=outputstyle]
Enter an integer: 45
45 is a positive number.
\end{lstlisting}

\finishExperiment

% ==========================================
% PROGRAM 80
% ==========================================

\startexperiment{20/08/2026}{WAP to accept any two integers from the keyboard and display the larger number.}{80.java}

\begin{lstlisting}[style=outputstyle]
Enter first integer: 55
Enter second integer: 99
Larger number = 99
\end{lstlisting}

\finishExperiment

% ==========================================
% PROGRAM 81
% ==========================================

\startexperiment{20/08/2026}{WAP to accept any integer from the keyboard and check whether it is odd or even.}{81.java}

\begin{lstlisting}[style=outputstyle]
Enter an integer: 36
36 is an even number.
\end{lstlisting}

\finishExperiment

% ==========================================
% PROGRAM 82
% ==========================================

\startexperiment{20/08/2026}{WAP to accept three numbers from the keyboard and display the largest number.}{82.java}

\begin{lstlisting}[style=outputstyle]
Enter first number: 33
Enter second number: 45
Enter third number: 99
Largest number = 99
\end{lstlisting}

\finishExperiment

% ==========================================
% PROGRAM 83
% ==========================================

\startexperiment{20/08/2026}{WAP to accept a number from the keyboard and display its multiplication table up to 10.}{83.java}

\begin{lstlisting}[style=outputstyle]
Enter a number: 20
20 x 1 = 20
20 x 2 = 40
20 x 3 = 60
20 x 4 = 80
20 x 5 = 100
20 x 6 = 120
20 x 7 = 140
20 x 8 = 160
20 x 9 = 180
20 x 10 = 200
\end{lstlisting}

\finishExperiment

% ==========================================
% PROGRAM 84
% ==========================================

\startexperiment{20/08/2026}{WAP to accept an integer from the keyboard and find and display the sum of its digits.}{84.java}

\begin{lstlisting}[style=outputstyle]
Enter an integer: 123
Sum of digits = 6
\end{lstlisting}

\finishExperiment

% ==========================================
% PROGRAM 85
% ==========================================

\startexperiment{20/08/2026}{WAP to accept an integer from the keyboard and find and display its factorial.}{85.java}

\begin{lstlisting}[style=outputstyle]
Enter an integer: 5
Factorial = 120
\end{lstlisting}

\finishExperiment

% ==========================================
% PROGRAM 86
% ==========================================

\startexperiment{20/08/2026}{WAP to create a 1D array, accept its elements from the keyboard, and find and display the largest element.}{86.java}

\begin{lstlisting}[style=outputstyle]
Enter size of array: 5
Enter element 1: 1
Enter element 2: 2
Enter element 3: 3
Enter element 4: 4
Enter element 5: 5
Largest element = 5
\end{lstlisting}

\finishExperiment

% ==========================================
% PROGRAM 87
% ==========================================

\startexperiment{20/08/2026}{WAP to create a 1D array, accept its elements from the keyboard, and find and display the smallest element.}{87.java}

\begin{lstlisting}[style=outputstyle]
Enter size of array: 5
Enter element 1: 5
Enter element 2: 4
Enter element 3: 3
Enter element 4: 2
Enter element 5: 1
Smallest element = 1
\end{lstlisting}

\finishExperiment

% ==========================================
% PROGRAM 88
% ==========================================

\startexperiment{20/08/2026}{WAP to create a 1D array, accept its elements from the keyboard, and find and display the sum and average of all elements.}{88.java}

\begin{lstlisting}[style=outputstyle]
Enter size of array: 5
Enter element 1: 1
Enter element 2: 2
Enter element 3: 3
Enter element 4: 4
Enter element 5: 5
Sum = 15
Average = 3.0
\end{lstlisting}

\finishExperiment

\end{document}